\begin{figure}[ht]
\centering
\begin{tikzpicture}[>=Stealth, font=\small, node distance=2.2cm]
  \node[draw, rounded corners, fill=blue!8, minimum width=4.6cm, minimum height=1cm] (riemann)
    {紧黎曼面 $X_0(N)(\C)$};
  \node[draw, rounded corners, fill=green!8, right=3.2cm of riemann, minimum width=4.8cm, minimum height=1cm] (curve)
    {复射影代数曲线 $C_{\C}\subset \mathbb P^n$};
  \draw[<->, thick] (riemann) -- node[above] {GAGA 视角} (curve);

  \node[draw, rounded corners, fill=orange!10, below=1.8cm of riemann, minimum width=4.8cm, minimum height=1cm] (qmodel)
    {$\Q$-模型 $C/\Q$};
  \node[draw, rounded corners, fill=purple!10, right=3.2cm of qmodel, minimum width=4.8cm, minimum height=1cm] (arith)
    {$C(\Q)$、$\Gal(\bar\Q/\Q)$、Hecke};

  \draw[->, thick] (curve) -- node[right] {系数域下降} (arith);
  \draw[->, thick] (qmodel) -- node[above] {复化} (curve);
  \draw[->, thick] (qmodel) -- node[above] {有理点与 Galois 作用} (arith);
\end{tikzpicture}
\caption{本章的核心转变：$X_0(N)$ 先作为复解析对象出现，但真正的算术问题要求它来自某条定义在 $\Q$ 上的代数曲线。只有在 $\Q$-模型上，Galois 作用与有理点才有内在意义。}
\label{fig:ch07-motivation}
\end{figure}
